% Draft repair, 24 November 2026; revised after one independent review; the revision itself is not reviewed
% !TEX root = ../monograph.tex
\subsection{Lattice symmetries and the average over forests}
\label{ch09:forest-average}

The expansion just described is carried out with one fixed forest,
that is, one rooted spanning tree in every block
(Section~\ref{ch05:actual-coordinates} and
Chapter~\ref{ch06:large-fields}). The block integral
$\mathscr Z_{\rm blk}$ does not involve a forest, but its individual
polymer activities do: the tree decides which fine links are
integration variables. A rooted tree also singles out directions in its
block. Indeed, for $L\ge3$ no spanning tree of the centered block which
is rooted at the center is invariant under the coordinate permutations:
the site $(1,1,1,1)$ is fixed by all of them, so its parent would be
fixed too, and none of its neighbors is. There is therefore no reason
for the lattice symmetries to transport the terms of a fixed-forest
expansion into one another.
Later in this chapter these symmetries are used term by term, namely in
the localization \eqref{ch09:localization} and in the identification of
the inverse coupling correction in Section~\ref{ch09:subtraction}.
We therefore average the polymer activities over the forests before
taking the logarithm. The average must be taken over independent
choices in the different blocks; an average over one global choice
would not preserve the product structure of the gas. This subsection
states what is averaged, which properties of the fixed-forest
construction are used, and what follows. One of those properties is
proved here, namely that the normalization constant does not depend on the
forest. The others are collected in
Assumption~\ref{ch09:forest-assumption}; they concern the construction
of Chapter~\ref{ch06:large-fields} and are not derived in this draft.
The paragraph at the end of the subsection lists what is proved, what
is assumed and what has not been adapted.

\subsubsection*{Symmetries and equivariant families}

Let $\Lambda$ be the coarse torus and $\Lambda^{\rm f}$ the fine torus,
and write $\mathbb B(y)=Ly+\mathbb B_{(L-1)/2}$ for the block of
$y\in\Lambda$. As in Chapter~\ref{ch01:lattice}, an element
$\varrho\in H_4$ and a site $z\in\Lambda$ define the bijection
$\varsigma(y)=\varrho y+z$ of $\Lambda$. We write $\varrho$ and
$\varsigma$ for these maps, instead of the letters $h$ and $g$, because
$h_X$ denotes the small-field indicator and $g$ a coupling in this
chapter. The maps $\varsigma$ form a finite group
$\mathfrak G=\mathfrak G(\Lambda)$; $\varsigma_z$ denotes the
translation $y\mapsto y+z$. The element $\varsigma$ also acts on fine
sites, by $x\mapsto\varrho x+Lz$; we use the same letter. Since $L$ is
odd and the blocks are centered, $\varsigma$ maps $\mathbb B(y)$ onto
$\mathbb B(\varsigma y)$, its center to the center, and internal block
edges to internal block edges. For a link configuration $U$ and an
ordered pair $(x,x')$ of nearest neighbors write $U(x,x')$ for the link
variable from $x$ to $x'$, so that $U(x',x)=U(x,x')^{-1}$. As in
Chapter~\ref{ch01:lattice}, $\varsigma$ acts on configurations by
\[
 (\varsigma^*U)(x,x')=U(\varsigma x,\varsigma x').
\]
Then $(\varsigma_1\varsigma_2)^*=\varsigma_2^*\varsigma_1^*$, a
function local on $X$ composed with $\varsigma^*$ is local on
$\varsigma X$, and
\begin{equation}
 \operatorname{Bad}_m(\varsigma^*U)
     =\varsigma^{-1}\operatorname{Bad}_m(U),
 \qquad h_X(\varsigma^*U)=h_{\varsigma X}(U).
 \label{ch09:bad-transport}
\end{equation}
The same definitions apply on the fine torus, and on the auxiliary tori
$\Lambda_i$ used below, each with its own group.

\begin{definition}[Equivariant families]
\label{ch09:equivariant}
A family of localized terms $(X_i,o_i,E_i)_{i\in I}$ is
\emph{equivariant} if for every $\varsigma\in\mathfrak G$ there is a
bijection $i\mapsto\varsigma\cdot i$ of $I$ such that
\[
 X_{\varsigma\cdot i}=\varsigma X_i,\qquad
 o_{\varsigma\cdot i}=\varsigma o_i,\qquad
 E_{\varsigma\cdot i}(U)=E_i(\varsigma^*U)\quad\text{for every }U .
\]
A family of large-field activities
$(Y_i,I_i^{\rm bad},\mathcal E_i^{\rm extra},K_i)_{i\in I}$ is
equivariant if there are such bijections with
$Y_{\varsigma\cdot i}=\varsigma Y_i$,
$I^{\rm bad}_{\varsigma\cdot i}=\varsigma I^{\rm bad}_i$,
$K_{\varsigma\cdot i}(U)=K_i(\varsigma^*U)$, and
$U\in\mathcal E^{\rm extra}_{\varsigma\cdot i}$ exactly when
$\varsigma^*U\in\mathcal E^{\rm extra}_i$. Representation data
are equivariant if both their families are. The same terminology is used
for any family of functions labelled with supports, and with base points
when they have them.
\end{definition}

Three elementary remarks will be used. First, a subfamily of an
equivariant family is equivariant when its index set is mapped to
itself by the bijections. Second, a fine family which is equivariant
for the group of the fine torus is in particular transported by the
maps $x\mapsto\varrho x+Lz$; only these enter a block step. Third, the
definition constrains the choice of base points. A support may be mapped to
itself by a symmetry that moves each of its sites, so one base point per
support cannot always be chosen. A function $f$ on such a support $X$ is
then divided into the $|X|$ equal shares $f/|X|$, one with base point at each
site of $X$, as is done for the selected terms in
Section~\ref{ch09:improvement}.

\subsubsection*{The normalization constant does not depend on the forest}

A forest $F$ consists of one spanning tree $F(y)$ of each block
$\mathbb B(y)$, made of internal block edges and rooted at $Ly$. Write
$\mathcal T_y$ for the finite set of such trees and
$\mathfrak T=\prod_{y\in\Lambda}\mathcal T_y$ for the set of forests.
As in Section~\ref{ch05:actual-coordinates}, $\iota_F$ extends
coordinates on the edges outside $F$ by zero on the tree edges.
Let $N_{\rm free}(F)$ be the number of positive fine edges outside $F$.

\begin{lemma}[Flat reference determinant]
\label{ch09:forest-determinant}
Let $\mathsf A^0$ be a real symmetric matrix acting on $\gLie$-valued
functions of the positive fine edges, and suppose that
$\mathsf A^0d_{0,f}\phi=0$ for every site function $\phi$ which vanishes
at all block roots. Then $N_{\rm free}(F)$ and
$\det(\iota_F^*\mathsf A^0\iota_F)$ do not depend on $F\in\mathfrak T$.
\end{lemma}

The hypothesis holds for the restriction of the flat quadratic form
\eqref{ch05:actual-hessian} to the fine variables. Write
$\mathcal Q_{\rm av}$ for the derivative of the averaging operation, which is
denoted $\mathcal Q$ in \eqref{ch05:actual-Q}; the subscript is added
because $\mathcal Q$ is the quadratic part of the path action in this
chapter. The form is
$d_{1,f}^*H_fd_{1,f}+b\,\mathcal Q_{\rm av}^*\mathcal Q_{\rm av}$,
with $b$ the coefficient of \eqref{ch05:actual-hessian}. A gradient has
no curl, and
$\mathcal Q_{\rm av}d_{0,f}\phi=d_{0,c}(\phi|_{L\Lambda})=0$ by
\eqref{ch05:actual-Q}. The hypothesis persists when further flat components
with the same property are added, and for each signed orientation,
because the elements of $H_4$ map gradients of functions vanishing at
the block roots to such gradients.

\begin{proof}
A spanning tree of a block has $L^4-1$ edges. Hence every forest has
$|\Lambda|(L^4-1)$ tree edges, and
$N_{\rm free}(F)=4L^4|\Lambda|-|\Lambda|(L^4-1)$.

Let $P_F$ be the coordinate projection onto functions supported on the
edges of $F$, and $D_F=1-P_F=\iota_F\iota_F^*$. For an edge function $u$
and a fine site $x$ let $\alpha_Fu(x)$ be the sum of $u$ along the path
in $F$ from $x$ to the root of its block, an edge traversed against its
orientation being counted with the opposite sign. With the convention
$(d_{0,f}\phi)(e)=\phi(e_+)-\phi(e_-)$, put
\[
 B_Fu=u+d_{0,f}\alpha_Fu .
\]
Summation along tree paths gives four facts: $\alpha_Fu$ vanishes at
every block root; $B_Fu$ vanishes on the edges of $F$; $B_Fu=u$ if $u$
already vanishes there, because then $\alpha_Fu=0$; and
$B_Fd_{0,f}\phi=0$ if $\phi$ vanishes at every block root, because then
$\alpha_Fd_{0,f}\phi=-\phi$.

Let $F,H\in\mathfrak T$. They have equally many tree edges; choose a
bijection between the edges of $F$ and those of $H$, and let $\Xi$
carry a function supported on the edges of $F$ to the function on the
edges of $H$ with the corresponding values. Define
\[
 T_{FH}u=\Xi P_Fu+B_HD_Fu .
\]
Since $B_HD_Fu$ vanishes on the edges of $H$, we have
$P_HT_{FH}u=\Xi P_Fu$ and $D_HT_{FH}u=B_HD_Fu$. Define $T_{HF}$ in
the same way with $\Xi^{-1}$. Then
\[
 T_{HF}T_{FH}u=P_Fu+B_FB_HD_Fu
   =P_Fu+B_F\bigl(D_Fu+d_{0,f}\alpha_HD_Fu\bigr)=P_Fu+D_Fu=u .
\]
The third equality uses the last two facts above, for the forest $F$,
and the vanishing of $\alpha_HD_Fu$ at the block roots. Exchanging the
roles of $F$ and $H$ gives $T_{FH}T_{HF}=1$. Both maps act on the edge
index by matrices with integer entries and trivially on $\gLie$. Their
determinants are therefore integers with product one, and
$|\det T_{FH}|=1$.

By hypothesis $\mathsf A^0B_HD_F=\mathsf A^0D_F$, and by symmetry of
$\mathsf A^0$ also $(B_HD_F)^*\mathsf A^0=D_F\mathsf A^0$. Moreover
$\langle T_{FH}u,P_HT_{FH}u\rangle=|\Xi P_Fu|^2=|P_Fu|^2$, so that
$T_{FH}^*P_HT_{FH}=P_F$. Hence
\[
 T_{FH}^*\bigl(D_H\mathsf A^0D_H+P_H\bigr)T_{FH}
   =D_F\mathsf A^0D_F+P_F .
\]
The matrix $D_F\mathsf A^0D_F+P_F$ is the direct sum of
$\iota_F^*\mathsf A^0\iota_F$ on the edges outside $F$ and the identity
on the edges of $F$. Taking determinants gives
$\det(\iota_F^*\mathsf A^0\iota_F)=\det(\iota_H^*\mathsf A^0\iota_H)$.
\end{proof}

The lemma is unconditional. Its application to the flat reference
determinant of this section is not: that determinant is described at
the beginning of the section in words, not by a formula, and that it is
$\det(\iota_F^*\mathsf A^0\iota_F)$ for a matrix of the kind just
described is part~(a) of the assumption below.

\subsubsection*{The fixed-forest expansion as a function of the forest}

For every $\varsigma\in\mathfrak G$ and every $y$, the map $\varsigma$
carries the trees of $\mathbb B(y)$ bijectively onto those of
$\mathbb B(\varsigma y)$. Thus
$(\varsigma F)(\varsigma y)=\varsigma(F(y))$ defines a bijection
$F\mapsto\varsigma F$ of $\mathfrak T$.

The Mayer expansion of this section was described for one forest. The
following assumption records, in the form of identities, its output,
its dependence on the forest, and the further properties of the
construction which the remainder of this subsection uses. In it,
$a_M^F$ denotes the integrated activity of a connected component $M$
computed with the forest $F$, summed over the remaining sector data as
described above. A \emph{local identification} between two tori
$\Lambda'$ and $\Lambda''$ is a bijection $\jmath$ between a set
$S'\subset\Lambda'$ and a set $S''\subset\Lambda''$, each contained in a
cube of side smaller than the period of its torus, which is the
restriction of a translation between these two cubes. It carries the
blocks of $S'$, their fine sites and links, their trees, and the link
variables read there to the corresponding objects for $S''$.

\begin{assumption}[Fixed-forest expansion]
\label{ch09:forest-assumption}
\mgubi{Issue 2: not proved. In existence/: (a) derived for one forest, its mask form asserted; (b), (e) argued; (c), (d) asserted; (f) a hypothesis. Lean has counterparts of (b), (c), (d), but for its own construction, not that of Chapter~\ref{ch06:large-fields}.}%
Under the hypotheses of Theorem~\ref{ch09:step-theorem} other than the
present assumption, with an equivariant incoming representation
data, the construction of this section has the following six
properties.
\begin{enumerate}[label=\textup{(\alph*)}]
\item \emph{Form.} For every forest $F\in\mathfrak T$ there are a
 countable set $\mathcal M$ of labels, each with a nonempty support
 $X_M\subset\Lambda$ containing its full dependence set and a
 prescribed large-field set $I_M\subset X_M$, possibly empty,
 all independent of $F$; real Borel gauge-invariant functions $a_M^F$
 of the coarse field, with $a_M^F$ local on $X_M$ and
 $\sum_M\sup_U|a_M^F(U)|<\infty$; a Borel function $\mathscr S$
 independent of $F$; and a constant $\mathscr N_F>0$, such that at every
 coarse field
 \begin{equation}
  \mathscr Z_{\rm blk}(U)=\mathscr N_Fe^{-\mathscr S(U)}
   \sum_{\mathcal A}\prod_{M\in\mathcal A}a_M^F(U).
  \label{ch09:fixed-forest-gas}
 \end{equation}
 The sum is over the finite sets $\mathcal A\subset\mathcal M$ whose
 supports are pairwise disjoint and for which
 $\bigcup_{M\in\mathcal A}I_M=\operatorname{Bad}_m(U)$. The constant is
 that of the bookkeeping at the beginning of this section:
 $\mathscr N_F=\mathscr N_0\,\mathfrak m^{N_{\rm free}(F)}
 \det(\iota_F^*\mathsf A^0\iota_F)^{-1/2}$, where $\mathscr N_0$ does
 not depend on $F$ and $\mathsf A^0$ is a matrix independent of $F$
 which satisfies the hypothesis of Lemma~\ref{ch09:forest-determinant}.
 For a label with $I_M=\varnothing$ there is a real Borel
 gauge-invariant function $f_M^F$, local on $X_M$, with
 \begin{equation}
  a_M^F=h_{X_M}f_M^F .
  \label{ch09:mask-form}
 \end{equation}
 Parts \textup{(b)}--\textup{(d)} hold for the functions $f_M^F$ as
 well as for the $a_M^F$.
\item \emph{Locality in the forest.} If $F(y)=F'(y)$ for every
 $y\in X_M$, then $a_M^F=a_M^{F'}$.
\item \emph{Uniform bounds.} The estimate \eqref{ch09:joint-input}, the
 majorants of the table following it, and the comparison estimates of
 Section~\ref{ch09:improvement} hold for every $F$ with right sides
 that do not depend on $F$.
\item \emph{Transport.} For every $\varsigma\in\mathfrak G$ there is a
 bijection $M\mapsto\varsigma\cdot M$ of $\mathcal M$ with
 \[
  X_{\varsigma\cdot M}=\varsigma X_M,\qquad
  I_{\varsigma\cdot M}=\varsigma I_M,\qquad
  a^{\varsigma F}_{\varsigma\cdot M}(U)=a^F_M(\varsigma^*U),
 \]
 and $\mathscr S(\varsigma^*U)=\mathscr S(U)$. The classification of
 labels used below (the number and kind of elementary factors, the
 presence of an incoming activity, and path lengths) is preserved by
 these bijections. A label which contains exactly one incoming localized
 term records its index $i$; the label $\varsigma\cdot M$ then
 records the index $\varsigma\cdot i$ given by the equivariance of the
 incoming family, and the coarse site $y$ with $o_i\in\mathbb B(y)$
 belongs to $X_M$.
\item \emph{Selected terms and agreement between tori.} The pure
 small-field construction, without incoming activities, is carried out
 with the same choices on each auxiliary torus $\Lambda_i$,
 $0\le i\le4$, of Section~\ref{ch09:improvement}, and parts
 \textup{(a)}--\textup{(d)} hold there with the group
 $\mathfrak G(\Lambda_i)$. On each torus $\Lambda'$ among
 $\Lambda,\Lambda_0,\ldots,\Lambda_4$ the selected low-order terms form
 a finite set. Its elements are tuples of labels with empty prescribed
 sets, and path terms which are terms of a path action averaged over
 $H_4$ as in Lemma~\ref{ch09:path-channels}. This set is mapped to
 itself by $\mathfrak G(\Lambda')$. For every local identification
 $\jmath$ between two of these tori, the labels occurring in selected
 tuples supported in $S'$ correspond bijectively to those supported in
 $S''$, with supports mapped by $\jmath$ and the classification of
 part~\textup{(d)} preserved; a tuple is selected exactly when the
 corresponding tuple is; and
 $f^{F''}_{\jmath M}(U'')=f^{F'}_M(U')$ whenever $U''$ corresponds to
 $U'$ and $F''(\jmath y)$ corresponds to $F'(y)$ for every $y\in X_M$.
 The selected path terms supported in $S'$ and in $S''$ correspond in
 the same way, with base points mapped by $\jmath$ and equal values at
 corresponding fields.
\item \emph{Preliminary families.} The regular functions $f_i$ of
 \eqref{ch09:early} are the terms of the logarithm of the averaged
 regular gas other than the selected ones, the tail of the truncation
 included, with the supports and base points of
 Proposition~\ref{ch09:family-equivariance}\textup{(ii)}, together
 with the localized terms $G_x$ of \eqref{ch09:localization} after the
 symmetrization of Lemma~\ref{ch09:symmetric-localization}. The preliminary large-field activities
 $(Y_i,I_i^{\rm bad},\mathcal E_i^{\rm extra},k_i)$ of
 \eqref{ch09:early} form an equivariant family.
\end{enumerate}
\end{assumption}

We comment on the six parts. Part~(a) is the reassembly described at
the beginning of this section and at the end of
Chapter~\ref{ch06:large-fields}; it is not displayed there as a formula,
and the objects $\mathscr S$, $\mathscr N_0$, $\mathsf A^0$ and the
labels are not identified there by name.
A label records the positions and kinds of its elementary factors. The
label set is made independent of $F$ by giving a factor which is absent
for a particular forest, such as a Haar remainder at a link belonging to
that forest, the value zero for that forest. The form
\eqref{ch09:mask-form} is the one used in the definition of $F_x$ and
$M_x$ in Section~\ref{ch09:improvement}, where the selected terms
appear both with and without their small-field indicator. The function
$f_M^F$ is an extension of $a_M^F$ beyond the configurations which are
small-field on $X_M$. Such an extension is a choice, and its properties
do not follow from those of $a_M^F$; this is why
\textup{(b)}--\textup{(d)} are required of it separately.
Part~(b) expresses that a
factor reads a tree only in the blocks which meet its dependence set.
For the Haar factors this requires comparing two local coordinate
systems as images of the same product Haar measure. Part~(c) is how the
constants of Chapters~\ref{ch05:gaussian} and~\ref{ch06:large-fields}
are stated: there are finitely many tree shapes at fixed $L$, and the
constants in \eqref{ch05:tree-comparison} are uniform in the allowed
trees. Part~(d) holds when every ingredient is transported together
with the forest. For the block kernel this is
Proposition~\ref{ch04:kernel-properties}; the path action
\eqref{ch04:path-action} is invariant; the incoming families are
equivariant by hypothesis. The remaining choices of
Chapter~\ref{ch06:large-fields}, namely the cutoff functions, the
finite-range covariances, the assignment of index triples to components
and the labels, must be made by rules which commute with $\mathfrak G$.
Part~(e) makes precise the phrase ``the same local formula on
$\Lambda_i$'' of Section~\ref{ch09:improvement}. The fixed-forest
agreement which it asserts rests on the fact that every quantity
entering a selected term is computed from data in its dependence
set: local windows of the quadratic form, their inverses, values
of the averaging operation and link words. This chapter and
Chapter~\ref{ch06:large-fields} contain no statement comparing two
tori, so the agreement is assumed. Part~(f) is needed because the
regrouping which leads from the gas to \eqref{ch09:early} is not
displayed as a formula in this chapter. For the large-field family it
is supported by part~(d) for the labels with $I_M\ne\varnothing$ and by
the fact that the deletion identity \eqref{ch08:deletion-identity}
refers only to supports.
We have not verified parts~(b)--(f) against every step of
Chapter~\ref{ch06:large-fields}; this is why they are stated as an
assumption.

\subsubsection*{The average}

Give $\mathfrak T$ the uniform probability measure. It is the product
of the uniform measures on the sets $\mathcal T_y$: the trees in
different blocks are chosen independently. Write $\mathbb E_F$ for the
corresponding average and set
\begin{equation}
 \bar a_M(U)=\mathbb E_Fa_M^F(U)
   =\frac1{|\mathfrak T|}\sum_{F\in\mathfrak T}a_M^F(U),
 \qquad
 \bar f_M(U)=\mathbb E_Ff_M^F(U)\quad\text{if }I_M=\varnothing .
 \label{ch09:averaged-activity}
\end{equation}

\begin{proposition}[Average over forests]
\label{ch09:forest-average-prop}
Assume parts \textup{(a)} and \textup{(b)} of
Assumption~\ref{ch09:forest-assumption}. Then:
\begin{enumerate}[label=\textup{(\roman*)}]
\item $\mathscr N_F=\mathscr N$ does not depend on $F$, and at every
 coarse field
 \begin{equation}
  \mathscr Z_{\rm blk}(U)=\mathscr Ne^{-\mathscr S(U)}
   \sum_{\mathcal A}\prod_{M\in\mathcal A}\bar a_M(U),
  \label{ch09:averaged-gas}
 \end{equation}
 with the same sets $\mathcal A$ as in \eqref{ch09:fixed-forest-gas}.
\item $\bar a_M$ is a real Borel gauge-invariant function local on
 $X_M$. It equals the average of $a_M^F$ over the trees $F(y)$,
 $y\in X_M$, alone. Every bound on $a_M^F$ or on its derivatives
 whose right side does not depend on $F$ holds for $\bar a_M$. Under
 part~\textup{(c)} this includes \eqref{ch09:joint-input}, with
 $\int G_M\,d\nu$ understood for each forest. If $I_M=\varnothing$,
 the same statements hold for $\bar f_M$, and
 $\bar a_M=h_{X_M}\bar f_M$.
\item Under part~\textup{(d)},
 $\bar a_{\varsigma\cdot M}(U)=\bar a_M(\varsigma^*U)$ for every
 $\varsigma\in\mathfrak G$, and likewise for $\bar f_M$ if
 $I_M=\varnothing$.
\end{enumerate}
\end{proposition}

\begin{proof}
(i) The first statement is Lemma~\ref{ch09:forest-determinant} applied
to the formula for $\mathscr N_F$ in part~(a). Fix $U$ and a finite set
$\mathcal A$ with pairwise disjoint supports. By part~(b), $a_M^F(U)$
is a function of the trees $(F(y))_{y\in X_M}$. For different
$M\in\mathcal A$ these are disjoint sets of coordinates of a product
measure, so
\[
 \mathbb E_F\prod_{M\in\mathcal A}a_M^F(U)
       =\prod_{M\in\mathcal A}\bar a_M(U).
\]
The index set of the sum in \eqref{ch09:fixed-forest-gas} does not
depend on $F$, and for each $F$ the sum converges absolutely, since
$\sum_{\mathcal A}\prod_{M\in\mathcal A}|a_M^F|
\le\exp(\sum_M\sup|a_M^F|)$. The left side of
\eqref{ch09:fixed-forest-gas} and the factor
$\mathscr Ne^{-\mathscr S(U)}$ do not depend on $F$. Averaging the
identity over the finite set $\mathfrak T$ and exchanging the finite
average with the absolutely convergent sum gives
\eqref{ch09:averaged-gas}.

(ii) A finite average of Borel, gauge-invariant, local functions has
these properties, and the triangle inequality gives the bounds. The
second statement is part~(b) and the product form of the measure. The
indicator $h_{X_M}$ does not depend on $F$, so averaging
\eqref{ch09:mask-form} gives $\bar a_M=h_{X_M}\bar f_M$; the other
statements about $\bar f_M$ follow as for $\bar a_M$, since
parts~(b) and~(c) are assumed for $f_M^F$.

(iii) The map $F\mapsto\varsigma F$ is a bijection of $\mathfrak T$,
so $\varsigma F$ is uniformly distributed when $F$ is. Therefore
\[
 \bar a_{\varsigma\cdot M}(U)
     =\mathbb E_Fa^{\varsigma F}_{\varsigma\cdot M}(U)
     =\mathbb E_Fa^F_M(\varsigma^*U)=\bar a_M(\varsigma^*U),
\]
and the same computation applies to $\bar f_M$.
\end{proof}

The product formula in the proof is where independence of the trees in
different blocks is used. It is not available for a general probability
measure on $\mathfrak T$, for instance for the uniform measure on the
translates and signed images of one reference forest.

From here on the polymer activities of the Mayer expansion are the averaged activities
$\bar a_M$. The identity \eqref{ch09:averaged-gas} has the form of
\eqref{ch09:fixed-forest-gas} with the same labels, supports, prescribed
sets, majorants and normalization constant. The cluster theorem and the
operations of the following sections therefore apply to it without
change. Two points deserve mention. An estimate of a single activity
which is made in one forest, with the same forest on both sides of a
comparison, passes to the average when it is uniform in the forest;
this is the case for the inherited term of
Section~\ref{ch09:improvement} under part~(c). By contrast, a cluster
term of the averaged gas with several polymers is a product of
independent averages and is not the average of the corresponding
fixed-forest cluster term. This does not affect
Sections~\ref{ch09:stationary-remainder}--\ref{ch09:centered-log}:
their estimate \eqref{ch09:pure-small-bound} concerns the logarithm of
the complete pure block integral, which involves no forest. Its
comparison with the sum of the selected terms is unchanged provided the
remaining terms are removed through their majorants, as in
Section~\ref{ch09:improvement}.

The second statement of (ii) also gives the comparison between tori
required for $F_y^{(i)}$ in \eqref{ch09:localization}. Let $\jmath$ be
a local identification as in part~(e) of the assumption, and let $M$ be
a label of a selected tuple supported in $S'$. The fixed-forest values
of $f_M$ and $f_{\jmath M}$ agree at corresponding fields when the
trees in the blocks of $X_M$ and of $\jmath X_M$ correspond. The
averaged functions $\bar f_M$ and $\bar f_{\jmath M}$ are averages over
the products of the sets of trees of these blocks with their uniform
measures, and $\jmath$ induces a bijection between the two products.
Hence $\bar f_{\jmath M}(U'')=\bar f_M(U')$ at corresponding fields.
Since $\jmath$ is injective, it also preserves the intersections of
supports, hence the incompatibility graph and the coefficient
\eqref{ch08:ursell} of a tuple, and the number of sites of a support,
hence the shares. The sums of selected terms at corresponding sites
therefore agree whenever all selected terms with a nonzero share at
these sites are supported in $S'$ and in $S''$.

\subsubsection*{Equivariance of the terms of the expansion}

The path action enters through its individual terms. Their transport
is elementary and is recorded first, since it is used twice below.

\begin{lemma}[Terms of the path action]
\label{ch09:path-channels}
Let the terms of $P^{\rm raw}$ in \eqref{ch04:path-action} be indexed
by $d$, with supports $Y_d$, base points $o_d$ and values
$L_d(U)$, as in Section~\ref{ch09:path-reset}. The terms of
$P^{\rm path}$ are then indexed by the pairs $(\varrho',d)$,
$\varrho'\in H_4$, with value $|H_4|^{-1}L_d(\varrho'^*U)$, support
$\varrho'Y_d$ and base point $\varrho'o_d$. This family is equivariant
under $\mathfrak G$. The same holds for the quadratic expressions
$a_dP_d$ without event indicators and for the cubic and quartic terms.
\end{lemma}

\begin{proof}
The expression $P^{\rm raw}$ is a sum over all sites of one expression
whose coefficients do not depend on the site. Hence the translate of a
term $d$ by $z$ is a term $d+z$, with $Y_{d+z}=Y_d+z$,
$o_{d+z}=o_d+z$ and $L_{d+z}(U)=L_d(\varsigma_z^*U)$. For
$\varrho''\in H_4$,
\[
 L_d(\varrho'^*\varrho''^*U)=L_d((\varrho''\varrho')^*U),
\]
so the value of $(\varrho',d)$ at $\varrho''^*U$ is the value of
$(\varrho''\varrho',d)$ at $U$, and
$(\varrho',d)\mapsto(\varrho''\varrho',d)$ is a bijection which maps
supports and base points by $\varrho''$. For a translation,
$\varsigma_z\varrho'=\varrho'\varsigma_{\varrho'^{-1}z}$ gives
$\varrho'^*\varsigma_z^*=\varsigma_{\varrho'^{-1}z}^*\varrho'^*$, so
the value of $(\varrho',d)$ at $\varsigma_z^*U$ is the value of
$(\varrho',d+\varrho'^{-1}z)$ at $U$. The support of the latter is
$\varrho'(Y_d+\varrho'^{-1}z)=\varrho'Y_d+z$, and similarly for the
base point. Translations and $H_4$ generate $\mathfrak G$. The argument
uses only the form of $P^{\rm raw}$ and applies to each of its parts.
\end{proof}

\begin{proposition}[Equivariant terms]
\label{ch09:family-equivariance}
Assume parts \textup{(a)}--\textup{(e)} of
Assumption~\ref{ch09:forest-assumption}, and form the gas from the
averaged activities.
\begin{enumerate}[label=\textup{(\roman*)}]
\item For a tuple $\Gamma=(M_1,\ldots,M_n)$ of labels with
 $I_{M_r}=\varnothing$, let
 $c_\Gamma=\frac1{n!}\varphi^T(\Gamma)\prod_r\bar a_{M_r}$ be its term
 in \eqref{ch08:log}, let
 $c^0_\Gamma=\frac1{n!}\varphi^T(\Gamma)\prod_r\bar f_{M_r}$, and let
 $X_\Gamma=\bigcup_rX_{M_r}$. Then $c_\Gamma=h_{X_\Gamma}c^0_\Gamma$.
 With
 $\varsigma\cdot\Gamma=(\varsigma\cdot M_1,\ldots,\varsigma\cdot M_n)$,
 \[
  X_{\varsigma\cdot\Gamma}=\varsigma X_\Gamma,\qquad
  c_{\varsigma\cdot\Gamma}(U)=c_\Gamma(\varsigma^*U),\qquad
  c^0_{\varsigma\cdot\Gamma}(U)=c^0_\Gamma(\varsigma^*U).
 \]
\item Let a set of such tuples be mapped to itself by every
 $\varsigma\in\mathfrak G$, and give each of its terms equal shares at
 the sites of its support, with base points there. The resulting family is
 equivariant. This applies to the terms of degree at least three and
 to the selected low-order cluster terms and their tail. The terms
 linear in an incoming localized term are not divided: such a term
 keeps the single coarse base point $y$ for which $\mathbb B(y)$ contains
 the base point of the incoming term, and with these base points they form
 an equivariant family as well. The selected path terms, with their
 base points, form an equivariant family. With $f_a^0$ running over the
 selected terms, namely the functions $c^0_\Gamma$ of the selected
 tuples and the selected path terms,
 \begin{equation}
  F_{\varsigma x}(U)=F_x(\varsigma^*U),\qquad
  M_{\varsigma x}(U)=M_x(\varsigma^*U)
  \label{ch09:share-equivariance}
 \end{equation}
 for the functions of Section~\ref{ch09:improvement}; in particular
 $F_x(1)$ does not depend on $x$.
\item The same statements hold on each auxiliary torus $\Lambda_i$
 with its group $\mathfrak G(\Lambda_i)$, for the pure small-field
 construction carried out there with its own average over forests.
\end{enumerate}
\end{proposition}

\begin{proof}
(i) If $X_\Gamma$ is small-field, every indicator $h_{X_{M_r}}$ equals
one. If it is not, a large-field site belongs to some $X_{M_r}$, and
the factor $\bar a_{M_r}=h_{X_{M_r}}\bar f_{M_r}$ vanishes. This gives
$c_\Gamma=h_{X_\Gamma}c^0_\Gamma$, also when labels are repeated. The
supports are transported by part~(d) of the assumption. Two labels are
incompatible when their supports intersect, and
$\varsigma X\cap\varsigma X'=\varsigma(X\cap X')$. Hence $\Gamma$ and
$\varsigma\cdot\Gamma$ have the same incompatibility graph and the
same coefficient \eqref{ch08:ursell}. The values are transported by
Proposition~\ref{ch09:forest-average-prop}(iii).

(ii) The bijection $\Gamma\mapsto\varsigma\cdot\Gamma$ preserves the
set by hypothesis, and maps the sites of $X_\Gamma$ onto those of
$X_{\varsigma\cdot\Gamma}$; the share $c_\Gamma/|X_\Gamma|$ with base point
$x$ is sent to the share of $\varsigma\cdot\Gamma$ with base point
$\varsigma x$. The set of tuples of degree at least three is mapped to
itself because part~(d) preserves the classification; the selected set
is mapped to itself by part~(e), and then so is its complement among
the low-order tuples, which forms the tail. For a term linear in an
incoming localized term with index $i$, part~(d) gives the index
$\varsigma\cdot i$ for the transported term. Its incoming base point is
$o_{\varsigma\cdot i}=\varsigma o_i$, which lies in
$\varsigma\mathbb B(y)=\mathbb B(\varsigma y)$ when $o_i\in\mathbb B(y)$.
So the coarse base point of the transported term is $\varsigma y$. The
selected path terms are terms of an averaged path action by
part~(e); they are permuted with their base points by
Lemma~\ref{ch09:path-channels}, and their selected set is mapped to
itself by part~(e). For \eqref{ch09:share-equivariance}, write
$a'=\varsigma\cdot a$; then $p_{\varsigma x,a'}=p_{xa}$,
$f^0_{a'}(U)=f^0_a(\varsigma^*U)$ by (i) and by the lemma, and
$h_{X_{a'}}(U)=h_{X_a}(\varsigma^*U)$ by \eqref{ch09:bad-transport};
the sum over selected $a'$ is a reindexing of the sum over selected
$a$. The flat configuration is fixed by $\varsigma^*$, and the
translations act transitively on $\Lambda$.

(iii) Part~(e) provides parts~(a)--(d) on $\Lambda_i$. The proofs of
(i) and (ii) and of Proposition~\ref{ch09:forest-average-prop} use no
property of the original torus that $\Lambda_i$ does not share.
\end{proof}

The base points in (ii) are those counted in
Section~\ref{ch09:improvement}: the bound for the inherited term there
counts the $L^4$ fine base points above one coarse base point, which is the
base point kept here. For the terms which are divided into equal shares
the following point is open. The norm $\|\cdot\|_m$ measures
derivatives from the base point, and placing a share of a cluster term at
an arbitrary site of its support, instead of at one distinguished
site, contributes at most a polynomial factor in the diameter of the support. That
this polynomial is absorbed by the margin in the support exponent in
\eqref{ch09:early-bounds} is used and has not been checked here.

The equivariance of the preliminary families of \eqref{ch09:early}
is now as follows. For the regular family $(f_i)$ it follows from
part~(f) of the assumption, from (ii), and from the symmetrization of
the localized terms carried out next. For the large-field family
$(k_i)$ it is part~(f) itself and is not proved here.

\subsubsection*{Symmetrizing the localization}

The functions $G_x$ and $L_x$ of \eqref{ch09:localization} are built
with the reflection map $q\mapsto x+\pi_{w+i}(q)$, whose image is the cube
$x+\{0,\ldots,w+i\}^4$. This cube is not symmetric about $x$: an axis
reflection about $x$ moves it by an amount which depends on $i$. The
family $(G_x)$ is transported by translations, but there is no reason
for it to be transported by axis reflections. Since the localization
identity is linear in $G_x$ and $L_x$, the missing symmetry is restored
by averaging.

\begin{lemma}[Symmetrized localization]
\label{ch09:symmetric-localization}
Suppose \eqref{ch09:share-equivariance} holds, and let $G_x,L_x$ be
defined by \eqref{ch09:localization}. Suppose, as is derived in
Section~\ref{ch09:improvement}, that
$L_x(U)=0$ whenever $N_{2\bar w}(x)$ is small-field and that, at every
$U$,
\begin{equation}
 \sum_xM_x(U)=|\Lambda|F_0(1)
   +\sum_xh_{N_{\bar w}(x)}(U)G_x(U)+\sum_xL_x(U).
 \label{ch09:localization-sum}
\end{equation}
For $\varrho\in H_4$ acting on $\Lambda$ by $y\mapsto\varrho y$, define
\[
 \widehat G_x(U)=\frac1{|H_4|}\sum_{\varrho\in H_4}
                        G_{\varrho^{-1}x}(\varrho^*U),
 \qquad
 \widehat L_x(U)=\frac1{|H_4|}\sum_{\varrho\in H_4}
                        L_{\varrho^{-1}x}(\varrho^*U).
\]
Then \eqref{ch09:localization-sum} holds with
$\widehat G_x,\widehat L_x$ in place of $G_x,L_x$;
$\widehat L_x(U)=0$ whenever $N_{2\bar w}(x)$ is small-field for $U$;
$\widehat G_x$ reads $U$ only on the cubes
$x+\varrho(\{0,\ldots,w+4\}^4)$, $\varrho\in H_4$; and the
families $(\widehat G_x)_x$, $(\widehat L_x)_x$, with base point $x$, are
equivariant under $\mathfrak G$.
\end{lemma}

\begin{proof}
Fix $\varrho$ and apply \eqref{ch09:localization-sum} to $\varrho^*U$.
By \eqref{ch09:share-equivariance},
$\sum_xM_x(\varrho^*U)=\sum_xM_{\varrho x}(U)=\sum_xM_x(U)$.
Neighborhoods are transported, $N_r(\varrho x)=\varrho N_r(x)$, so
$h_{N_{\bar w}(x)}(\varrho^*U)=h_{N_{\bar w}(\varrho x)}(U)$ by
\eqref{ch09:bad-transport}. Substituting $x=\varrho^{-1}x'$ in the two
site sums gives
\[
 \sum_xM_x(U)=|\Lambda|F_0(1)
  +\sum_{x'}h_{N_{\bar w}(x')}(U)G_{\varrho^{-1}x'}(\varrho^*U)
  +\sum_{x'}L_{\varrho^{-1}x'}(\varrho^*U).
\]
Averaging over $\varrho$ proves the identity. If $N_{2\bar w}(x)$ is
small-field for $U$, then $N_{2\bar w}(\varrho^{-1}x)$ is small-field
for $\varrho^*U$, so every term of $\widehat L_x(U)$ vanishes. The
function $U\mapsto G_{\varrho^{-1}x}(\varrho^*U)$ is the expression for
$G_x$ with the reflection map $q\mapsto x+\varrho\pi_{w+i}(q)$ in place of
$q\mapsto x+\pi_{w+i}(q)$, which gives the dependence set.

For $\varrho''\in H_4$, substituting $\varrho=\varrho''\varrho'$ gives
\[
 \widehat G_{\varrho''x}(U)=\frac1{|H_4|}\sum_{\varrho'}
        G_{\varrho'^{-1}x}(\varrho'^*\varrho''^*U)
                   =\widehat G_x(\varrho''^*U).
\]
By construction $G_{x+z}(U)=G_x(\varsigma_z^*U)$ for the translation
$\varsigma_z$, because $U_{i,x+z}=(\varsigma_z^*U)_{i,x}$.
Since $\varsigma_z\varrho=\varrho\varsigma_{\varrho^{-1}z}$,
\[
 G_{\varrho^{-1}(x+z)}(\varrho^*U)
   =G_{\varrho^{-1}x}(\varsigma_{\varrho^{-1}z}^*\varrho^*U)
   =G_{\varrho^{-1}x}(\varrho^*\varsigma_z^*U),
\]
and averaging gives
$\widehat G_{x+z}(U)=\widehat G_x(\varsigma_z^*U)$.
Translations and $H_4$ generate $\mathfrak G$. The same computation
applies to $L_x$, whose translation covariance follows from that of
$M_x$, of $F_x(1)$ and of the last sum in \eqref{ch09:localization}.
\end{proof}

The bounds stated for $G_x$ and $L_x$ in
Section~\ref{ch09:improvement} are used for $\widehat G_x$ and
$\widehat L_x$. For $\widehat L_x$ the supremum bound passes to the
average directly. For $\widehat G_x$ the reason is that the bound on
$G_x$ is obtained from the bound on the auxiliary sums and from the
reflection map, of which it uses only that nearest neighbors are mapped
to nearest neighbors within distance $w+i$ of the base point; the maps
$q\mapsto x+\varrho\pi_{w+i}(q)$ have this property. The derivation of
the bound on $G_x$ is given in Section~\ref{ch09:improvement} in one
sentence, and the corresponding statement for the maps with $\varrho$
has the same status: it is not proved in detail here. The
dependence set of $\widehat G_x$ extends to distance $w+4$ on all
sides of $x$; it is contained in the neighborhood $N_{\bar w}(x)$
already used for its indicator.

\subsubsection*{The summed affine Hessian}

For $\varrho\in H_4$ write $\varrho e_a=\omega_ae_{\varrho(a)}$; this
defines signs $\omega_a=\pm1$ and a permutation $a\mapsto\varrho(a)$ of
the four axes. For an orientation $q=(a,b)$, $a<b$, let $\varrho\cdot q$
be the pair $\{\varrho(a),\varrho(b)\}$ in increasing order, and
\[
 s_\varrho(q)=\omega_a\omega_b
        \operatorname{sgn}(\varrho(b)-\varrho(a)).
\]
With the convention $t_{ba}=-t_{ab}$ for an antisymmetric array $t$,
this means
$\omega_a\omega_bt_{\varrho(a)\varrho(b)}=s_\varrho(q)t_{\varrho\cdot q}$.

\begin{lemma}[Isotropy of the summed coefficient]
\label{ch09:isotropy}
Let $(X_i,o_i,f_i)_{i\in I}$ be an equivariant family of preliminary
regular functions as in \eqref{ch09:early}, with
$\sum_{o_i=o}|q_{f_i}(q,r)|<\infty$ for every $o$, $q$, $r$. Let
$\varsigma(y)=\varrho y+z$.
\begin{enumerate}[label=\textup{(\roman*)}]
\item $q_{f_{\varsigma\cdot i}}(\varrho\cdot q,\varrho\cdot r)
 =s_\varrho(q)s_\varrho(r)\,q_{f_i}(q,r)$, and
 $F^{\rm cl}_q(\varsigma^*U,o)
 =s_\varrho(q)F^{\rm cl}_{\varrho\cdot q}(U,\varsigma o)$.
\item There is a real number $\Delta$ with
 $\sum_{o_i=o}q_{f_i}(q,r)=\Delta\delta_{qr}$ for every $o$.
\end{enumerate}
Part \textup{(ii)} uses three kinds of symmetries: translations give
independence of $o$; the reflection of one axis about $o$ gives the
vanishing for $q\ne r$; coordinate permutations about $o$, without
signs, give equality of the diagonal entries.
\end{lemma}

\begin{proof}
(i) Put $Z_i=X_i\cup N_2(o_i)$; then
$Z_{\varsigma\cdot i}=\varsigma Z_i$, and one of them wraps exactly
when the other does, in which case both coefficients vanish by
definition. Otherwise the coordinate lifts are displacements in $\ZZ^4$
and $\xi_{\varsigma o}(\varsigma y)=\varrho\xi_o(y)$. Let $A$ be an
affine potential on $Z_{\varsigma\cdot i}$ with
$A(y,y+e_c)=u_c+\sum_\nu\xi_{\varsigma o}(y)^\nu v_{c\nu}$, and
$A(y',y)=-A(y,y')$. Its curl is the antisymmetric array
$(\mathrm dA)_{ab}=v_{ba}-v_{ab}$. The pullback
$(\varsigma^*A)(x,x')=A(\varsigma x,\varsigma x')$ satisfies
$\exp(\varsigma^*A)=\varsigma^*\exp A$. If $\omega_a=1$,
$(\varsigma^*A)(x,x+e_a)=A(\varsigma x,\varsigma x+e_{\varrho(a)})$;
if $\omega_a=-1$,
$(\varsigma^*A)(x,x+e_a)=-A(\varsigma x-e_{\varrho(a)},\varsigma x)$,
which differs from $-A(\varsigma x,\varsigma x+e_{\varrho(a)})$ by the
constant $v_{\varrho(a)\varrho(a)}$. In both cases, since
$(\varrho\xi)^{\varrho(\nu)}=\omega_\nu\xi^\nu$,
\[
 (\varsigma^*A)(x,x+e_a)=u'_a+\sum_\nu\xi_o(x)^\nu
     \omega_a\omega_\nu v_{\varrho(a)\varrho(\nu)} .
\]
So $\varsigma^*A$ is affine on $Z_i$ with curl
$(\mathrm d\varsigma^*A)_{ab}
=\omega_a\omega_b(\mathrm dA)_{\varrho(a)\varrho(b)}$, that is,
$(\mathrm d\varsigma^*A)_q=s_\varrho(q)(\mathrm dA)_{\varrho\cdot q}$.

By Section~\ref{ch09:subtraction}, $f$ and $Q_f$ have the same
Hessian on affine potentials. The derivative of $F^{\rm cl}_q$ at the
identity is the curl component $(\mathrm dA)_q$, and $q_f$ is
symmetric. Hence, for affine potentials $A,\widetilde A$,
\[
 D^2[f(\exp\cdot)](0)[A,\widetilde A]
     =\sum_{q,r}q_f(q,r)
      \langle(\mathrm dA)_q,(\mathrm d\widetilde A)_r\rangle .
\]
Apply this to $f_{\varsigma\cdot i}=f_i\circ\varsigma^*$:
\[
 \sum_{q',r'}q_{f_{\varsigma\cdot i}}(q',r')
       \langle(\mathrm dA)_{q'},(\mathrm d\widetilde A)_{r'}\rangle
 =\sum_{q,r}q_{f_i}(q,r)s_\varrho(q)s_\varrho(r)
       \langle(\mathrm dA)_{\varrho\cdot q},
              (\mathrm d\widetilde A)_{\varrho\cdot r}\rangle .
\]
The curls are arbitrary and $q\mapsto\varrho\cdot q$ is a bijection of
the six orientations, which gives the first formula.

For the clover curvature, the loop of $\varsigma^*U$ which starts at
$o$ and follows
$\varepsilon e_a,\eta e_b,-\varepsilon e_a,-\eta e_b$ is the loop of
$U$ which starts at $\varsigma o$ and follows
$\varepsilon\omega_ae_{\varrho(a)},\eta\omega_be_{\varrho(b)}$ and
their negatives. Substituting $\varepsilon'=\varepsilon\omega_a$,
$\eta'=\eta\omega_b$ in the definition gives the factor
$\omega_a\omega_b$. If $\varrho(a)>\varrho(b)$ the loop is the
inverse of the one in the definition of
$F^{\rm cl}_{\varrho(b)\varrho(a)}$, and
$\operatorname{proj}_{\gLie}(u-u^{-1})$ is odd under inversion. This is
the second formula.

(ii) Put $\Delta_o(q,r)=\sum_{o_i=o}q_{f_i}(q,r)$. The bijection
$i\mapsto\varsigma\cdot i$ maps $\{i:o_i=o\}$ onto
$\{i:o_i=\varsigma o\}$. Reindexing the absolutely convergent sum and
using (i),
\[
 \Delta_{\varsigma o}(\varrho\cdot q,\varrho\cdot r)
     =s_\varrho(q)s_\varrho(r)\Delta_o(q,r).
\]
For a translation, $\varrho=1$ and $\Delta_{o+z}=\Delta_o$. Let
$q\ne r$. They are distinct two-element subsets of $\{1,2,3,4\}$, so
some axis $c$ belongs to exactly one of them. Let $\varsigma$ be the
reflection of the axis $c$ about $o$. Then $\varsigma o=o$,
$\varrho\cdot q=q$, $\varrho\cdot r=r$ and
$s_\varrho(q)s_\varrho(r)=-1$, so $\Delta_o(q,r)=0$. Finally let
$q,q'$ be two orientations, and let $\varsigma$ fix $o$ and have as
linear part a permutation matrix $\varrho$, without signs, which
carries the pair $q$ to the pair $q'$. Then all $\omega_a=1$, so
$s_\varrho(q)^2=1$ and $\Delta_o(q',q')=\Delta_o(q,q)$.
\end{proof}

\subsubsection*{Equivariance of the output}

\begin{proposition}[The step preserves equivariance]
\label{ch09:output-equivariance}
Suppose the families $(X_i,o_i,f_i)$ and
$(Y_i,I_i^{\rm bad},\mathcal E_i^{\rm extra},k_i)$ of
\eqref{ch09:early} are equivariant, and let $L_x$ denote the
symmetrized corrections of Lemma~\ref{ch09:symmetric-localization}.
Then the localized terms and the large-field activities of the
output tuple constructed in
Sections~\ref{ch09:subtraction}--\ref{ch09:fusion} form equivariant
families.
\end{proposition}

\begin{proof}
By Lemma~\ref{ch09:isotropy}(i),
\[
 Q_{f_i}(\varsigma^*U)
   =\tfrac12\sum_{q,r}q_{f_i}(q,r)s_\varrho(q)s_\varrho(r)
   \langle F^{\rm cl}_{\varrho\cdot q}(U,\varsigma o_i),
           F^{\rm cl}_{\varrho\cdot r}(U,\varsigma o_i)\rangle
   =Q_{f_{\varsigma\cdot i}}(U).
\]
With \eqref{ch09:bad-transport} and
$f_{\varsigma\cdot i}(1)=f_i(1)$ this gives
$R_{f_{\varsigma\cdot i}}(U)=R_{f_i}(\varsigma^*U)$ and
$\zeta_{f_{\varsigma\cdot i}}(U)=\zeta_{f_i}(\varsigma^*U)$, on the
supports $Z_{\varsigma\cdot i}=\varsigma Z_i$.

By Lemma~\ref{ch09:path-channels}, $\mathfrak G$ permutes the terms
of the path action, with their supports and base points, and transports
the quadratic expression $a_dP_d$ together with its quadratic term. The clover
substitute of a quadratic term is transported by
Lemma~\ref{ch09:isotropy}(i). Since $\alpha$ and $\Delta$ are numbers,
$F_d$ and $J_d$ are transported as well, and likewise the cubic and
quartic terms.

Hence the large-field corrections $\ell_a$ of
Section~\ref{ch09:fusion}, which are the $\zeta_{f_i}$, the $J_d$ and
the symmetrized $L_x$, form an equivariant family with supports $Z_a$.
A finite connected family $P$ of preliminary large-field labels and
correction labels is mapped by $\varsigma$ to the family
$\varsigma\cdot P$ of the transported labels. It is again connected,
with $Y_{\varsigma\cdot P}=\varsigma Y_P$ and
$I_{\varsigma\cdot P}=\varsigma I_P$, and by
\eqref{ch09:bad-transport}
\[
 \mathbf1_{\{Y_{\varsigma\cdot P}\cap\operatorname{Bad}_m(U)
              =I_{\varsigma\cdot P}\}}
 =\mathbf1_{\{Y_P\cap\operatorname{Bad}_m(\varsigma^*U)=I_P\}} .
\]
The products in $K_P$ and the inherited additional conditions are
transported factor by factor. So
$K_{\varsigma\cdot P}(U)=K_P(\varsigma^*U)$. The
localized terms of the output are the functions
$h_{Z_i}R_{f_i}$ and the regular terms composing $F$, which
have been treated.
\end{proof}

At the Wilson start of Section~\ref{ch09:wilson} the regular term of a
plaquette $p$ is a class function of $U_p$ which is even under
inversion, so it depends neither on the base point nor on the
orientation of $p$. Divided into four equal shares with base points at the
vertices of $p$, these terms form an equivariant family. The components
$I$ of large-field sites, their supports $N_{2B}(I)$ and their
energies $d_I$ are defined by rules which commute with $\mathfrak G$,
so the family $(K_I)$ is equivariant. Under
Assumption~\ref{ch09:forest-assumption},
Propositions~\ref{ch09:family-equivariance}
and~\ref{ch09:output-equivariance} then propagate equivariance from
each scale index to the next.

\paragraph{Status of this subsection.}
\emph{Proved here, with complete proofs.}
Lemma~\ref{ch09:forest-determinant}: the number of free edges and the
flat determinant are the same for all forests.
Lemma~\ref{ch09:path-channels}: the lattice symmetries permute the
terms of the averaged path action.
Lemma~\ref{ch09:symmetric-localization}: the localization identity
survives the average of $G_x$ and $L_x$ over $H_4$, and the averaged
families are equivariant.
Lemma~\ref{ch09:isotropy}: for an equivariant family the summed
quadratic coefficient is a multiple of the identity.
Proposition~\ref{ch09:output-equivariance}: subtraction, path
replacement and exact regrouping preserve equivariance.
Propositions~\ref{ch09:forest-average-prop}
and~\ref{ch09:family-equivariance} are proved from
Assumption~\ref{ch09:forest-assumption} and are no stronger than it.

\emph{Assumed.} The six parts of
Assumption~\ref{ch09:forest-assumption}, none of which is proved in
this monograph at present.
Part~(a): the exact identity for the block integral with one forest,
the formula for its normalization constant, and the form
\eqref{ch09:mask-form} of the activities without large-field sites;
the beginning of this section and the last section of
Chapter~\ref{ch06:large-fields} describe this reassembly in words and
state that its identities are separate from the bounds proved there.
Part~(b): an activity reads the forest only in the blocks of its
support; Chapter~\ref{ch06:large-fields} defines the dependence set
of a factor as containing every window entry and frozen coordinate
which it reads, but does not state the consequence.
Part~(c): bounds uniform in the forest; Chapter~\ref{ch05:gaussian}
takes common constants over the finitely many tree shapes and states
\eqref{ch05:tree-comparison} uniformly in the allowed trees, while the
uniformity of \eqref{ch09:joint-input} and of the comparison estimates
is not stated.
Part~(d): transport of the activities by the lattice symmetries
together with the forest; proved for the block kernel
(Proposition~\ref{ch04:kernel-properties}) and built into the path
action \eqref{ch04:path-action}, not examined for the cutoffs,
covariances and assignment rules of Chapter~\ref{ch06:large-fields}.
Part~(e): the same construction on the auxiliary tori, selected sets
which are mapped to themselves and correspond between tori, and equal
fixed-forest values under local identifications; supported only by the
locality of the constructions of Chapters~\ref{ch05:gaussian}
and~\ref{ch06:large-fields}, which contain no comparison of two tori.
Part~(f): the list of the preliminary regular functions and the
equivariance of the preliminary large-field activities;
the regrouping which produces them is not displayed in this chapter.

\emph{Used and not proved, apart from the assumption.} The vanishing
of $L_x$ on small-field neighborhoods and the identity
\eqref{ch09:localization-sum}, which Section~\ref{ch09:improvement}
derives by a classification of sites that is only sketched there; the
derivative bound for the localized terms after the average over $H_4$;
and the absorption, by the margin in the support exponent, of the polynomial factor from placing the base points of
equal shares at all sites of a support.

\emph{Not yet adapted.}\mgubi{Out of date: Chapter~\ref{ch11:insertions} now averages over the forests, under Hypothesis~(H) of Section~\ref{ch11:marked-step} (not proved). The remark on Section~\ref{ch17:quadratic-current} stands (issue 7).} The expansions with insertions of
Chapter~\ref{ch11:insertions}, which fix one forest: their activities
with insertions must be averaged over the forests in the same way, so that
at zero source they reduce to the representation constructed here; no
such statement has been written. The histories of
Chapter~\ref{ch10:trajectory} depend on the assumption through
Theorem~\ref{ch09:step-theorem}. Chapter~\ref{ch06:large-fields} does
not yet say that its forest is averaged afterwards. The zero-background
computation of Section~\ref{ch17:quadratic-current} uses one translated
forest and relies on an unproved independence of the forest for
gauge-invariant quantities.
